% ============================================================
%  Appendix
% ============================================================
\chapter{แบบสอบถามการประเมินระบบ}

แบบสอบถามที่ใช้ในการประเมินระบบมีหัวข้อดังนี้:

\textbf{ส่วนที่ 1 — ข้อมูลทั่วไปของผู้ทดสอบ}
\begin{itemize}
  \item ตำแหน่ง/บทบาท
  \item ระดับประสบการณ์กับการอ่านผล AST (Disk Diffusion)
\end{itemize}

\textbf{ส่วนที่ 2 — การทดสอบตามสถานการณ์ (Scenario-Based Test)}
\begin{enumerate}
  \item การอัปโหลดภาพและสั่งรันระบบวิเคราะห์
  \item การแก้ไขผลลัพธ์ด้วยตนเอง (Manual Correction)
  \item การดูประวัติและดาวน์โหลดรายงาน
\end{enumerate}

\textbf{ส่วนที่ 3 — คะแนนความพึงพอใจ (Likert Scale 1--5)}
\begin{enumerate}
  \item ขั้นตอนการอัปโหลดและวิเคราะห์ทำได้ง่ายและไม่ซับซ้อน
  \item ตารางแสดงผลจัดวางได้เป็นระเบียบ อ่านข้อมูลได้ชัดเจน
  \item ฟังก์ชัน Manual Correction ใช้งานสะดวก
  \item ระบบช่วยลดเวลาและภาระงานเมื่อเทียบกับการวัดด้วยไม้บรรทัด
  \item มั่นใจในความถูกต้องของการคำนวณและการแปลผล
  \item โดยภาพรวมระบบมีประโยชน์และพร้อมใช้งานจริงในห้องปฏิบัติการ
\end{enumerate}

\textbf{ส่วนที่ 4 — ความคิดเห็นเพิ่มเติม}
\begin{itemize}
  \item จุดเด่น/สิ่งที่ชอบในระบบ
  \item ปัญหาที่พบ/สิ่งที่ทำให้หงุดหงิด
  \item ฟีเจอร์เพิ่มเติมที่ต้องการ
  \item ข้อเสนอแนะอื่นๆ
\end{itemize}
